% Section 6: Austin's own money for music. Neutral description of what the
% money did and when, with innocent explanations noted where they exist.
\section{Austin's music money is small, uneven, and has never had a performance audit}
\label{sec:city}

Austin funds music from its Hotel Occupancy Tax, a levy that lodging guests
pay on hotel and short-term-rental stays inside the city. A 2019 ordinance
raised the tax and directed a fixed share of the increase to a Live Music
Fund, so the money moves with hotel demand rather than with any annual
council judgment about musicians.\footnote{City of Austin Ordinance
No.~20190919-149, adopted 19 September 2019 and effective 30 September
2019, raised the Hotel Occupancy Tax rate from 7 to 9 percent and directed
an amount equal to 15 percent of the two-point increase to the Live Music
Fund, another 15 percent to Historic Preservation, and the balance largely
to Convention Center expansion. Because the ordinance fixes the share
rather than reappropriating each year, fund revenue tracks hotel demand.
Ordinance text at
\url{https://services.austintexas.gov/edims/document.cfm?id=328565}.}

\subsection{Live music receives a small share of a large tax}

\begin{figure}[htbp]
\centering
\includegraphics[width=\textwidth]{fig12_hot_allocation.png}
\caption{City of Austin Hotel Occupancy Tax revenue by destination, fiscal
year 2025, nominal dollars. Source: City of Austin Open Budget ATX,
\textit{Revenue Open Budget -- revised} (Socrata dataset 5rhv-xasu),
fund-level actuals for fiscal 2025. How to read it: all bars share one
zero-based axis, so the relative sizes are directly comparable. Caveat: the
residual bar is total hotel-tax revenue minus the three named funds rather
than a single program, and under the 2019 ordinance most of the residual
flows to Convention Center expansion. Because every bar is same-year,
nominal and real dollars give the same picture.}
\label{fig:hotalloc}
\end{figure}

In fiscal 2025 the Live Music Fund received \NCityLmfRevenueFyTwoFive, or
\NCityLmfShareOfHotFyTwoFive\ of the city's
\NCityHotTotalRevenueFyTwoFive\ in Hotel Occupancy Tax
revenue.\footnote{The city's Open Budget ATX portal posts fund-level
actuals for every named city fund by fiscal year; those posted actuals are
the auditable primary record until each year's Annual Comprehensive
Financial Report closes. Fund revenue and Hotel Occupancy Tax totals for
fiscal 2025 in this section come from that portal's \textit{Revenue Open
Budget -- revised} extract, Socrata dataset 5rhv-xasu, retrieved 1 August
2026 from \url{https://data.austintexas.gov/resource/5rhv-xasu.csv}. Same-
year ratios need no deflation because both sides carry the same dollar
base. The Open Budget series begins in fiscal 2021 and does not observe
any earlier fund balance.} The Historic Preservation Fund drew
\NCityHpRevenueFyTwoFive\ from the same tax, \NCityHpToLmfRatioFyTwoFive\
as much as the Live Music Fund. That gap is not a departure from the 2019
ordinance's equal split. The Historic Preservation Fund also receives
hotel-tax money beyond the 15 percent of the two-point increase the
ordinance dedicated to each fund, and on the equal share alone the two
would draw the same amount.

From fiscal 2024 to fiscal 2025, Live Music Fund revenue moved
\NCityLmfRevenuePctchgFyTwoFourFyTwoFiveNominal\ in nominal terms and
\NCityLmfRevenuePctchgFyTwoFourFyTwoFiveReal\ in 2025
dollars.\footnote{Fiscal-year change computed from the same Open Budget
revenue extract as above (dataset 5rhv-xasu). This section approximates
Austin fiscal year \(Y\) by calendar year \(Y\) for deflation, because
Austin fiscal years run October through September and nine of the twelve
months fall in calendar \(Y\); no quarterly deflator was available for the
national consumer price index used here.} That is one year on a series
that tracks hotel demand and has swung widely since 2021, so it should not
be read as a trend.

\subsection{The fund collected for two observed years before it paid musicians}

\begin{figure}[htbp]
\centering
\includegraphics[width=\textwidth]{fig13_lmf_flow.png}
\caption{Live Music Fund revenue against expenditure by fiscal year, 2025
dollars, with the cumulative position on the right axis. Sources: City of
Austin Open Budget ATX, \textit{Revenue Open Budget -- revised} (dataset
5rhv-xasu) and \textit{Expense Open Budget -- revised} (dataset
yeeq-kk6v), Live Music Fund rows for fiscal 2021 through fiscal 2026. How
to read it: bars are annual flows and the line is the running total of
revenue minus expenditure. Caveat: that cumulative line is a cash-flow
proxy built from fiscal 2021 forward, not an audited fund balance. It
excludes any opening balance and does not reconcile with press reporting
that put the fund's balance at roughly \$3.2 million in September 2025, a
gap that cannot be resolved without the city's audited fund-balance
schedule.}
\label{fig:lmfflow}
\end{figure}

The City Council created the Live Music Fund in September 2019 and the
fund did not disburse a substantial grant until fiscal 2023. By the close
of fiscal 2022 the fund had accumulated \NCityLmfAccumByFyTwoTwoNominal\
in unspent nominal money across the period when musicians' income was
most disrupted.\footnote{The Open Budget series begins in fiscal 2021, so
fiscal 2020 collections are not in that figure. The figure is a cumulative
revenue-minus-expenditure proxy built from the two Open Budget extracts
(5rhv-xasu and yeeq-kk6v) rather than an audited fund balance, and it
excludes any opening balance the fund carried from its 30 September 2019
creation. Figure~\ref{fig:lmfflow} plots the same cumulative position in
2025 dollars, where it reaches a higher number for the same fiscal year.
The gap against press reporting of a \$3.2 million balance in September
2025 remains open pending the city's Annual Comprehensive Financial
Report fund-balance schedule.} Innocent explanations for the delay exist:
a new fund has to write eligibility rules, procure an administrator and
stand up an application process, and the pandemic disrupted all three.
The timing still matters. The fund spent \$4,284 in fiscal 2021 and
\$50,000 in fiscal 2022 against \$1.5 million and \$3.7 million in
revenue,\footnote{Annual amounts in this sentence are from the Open Budget
expense (yeeq-kk6v) and revenue (5rhv-xasu) extracts,
\texttt{fund\_nm~=~Live~Music~Fund}, actuals for fiscal 2021 and fiscal
2022.} so almost nothing left the fund during the years musicians' income
was most disrupted.

Of every dollar the fund has spent through fiscal 2026, roughly 13 cents
paid for something other than grants to musicians, venues and promoters:
\NCityLmfAdminShareGrants\ of cumulative spending went to grants to
subrecipients, \NCityLmfAdminShareConsultant\ to the third-party
administrator that runs the application process, and
\NCityLmfAdminShareAdvertising\ to advertising.\footnote{Object-level
actuals for fiscal 2021 through fiscal 2027, summed across years, from the
City of Austin Open Budget ATX expense extract (dataset yeeq-kk6v) with
\texttt{fund\_nm~=~Live~Music~Fund}, retrieved 1 August 2026 from
\url{https://data.austintexas.gov/resource/yeeq-kk6v.csv}. The three
denominated objects are Grants to Subrecipients, Consultant-Others (the
fund's third-party administrator, the Long Center for the Performing
Arts, from fiscal 2023 onward) and Advertising. The Open Budget expense
series carries no Personnel or other object rows for this fund. Fiscal
2026 is a partial-year actual, unclosed at the time of writing, and
fiscal 2027 rows are budget only. Fiscal 2026 figures for the three
shares thus draw on a year that had not closed when this study was
written.}

\subsection{What a musician could receive depended on the year they applied}

\begin{figure}[htbp]
\centering
\includegraphics[width=\textwidth]{fig14_lmf_cohorts.png}
\caption{Live Music Fund awards and average award by cycle. Sources: the
city's fiscal 2023, fiscal 2024 and fiscal 2026 awardee lists, using each
document's own stated headline totals, plus the official Live Music Fund
program page for the fiscal 2025 zero. How to read it: bars are the
number of awards on the left axis and the line is the average award on
the right axis, broken rather than interpolated at fiscal 2025 because no
cycle ran that year. Caveats: the fiscal 2023 count is
\NCityLmfFyTwoThreeAwardsList\ in the awardee list PDF against
\NCityLmfFyTwoThreeAwardsPress\ in press reporting, a difference that
appears to be one duplicated name; the fiscal 2026 list does not
reconcile with itself, so its stated headline totals are used here.}
\label{fig:lmfcohorts}
\end{figure}

From one cycle to the next, the average award and the number of
recipients moved in opposite directions, and no cycle ran at all in
fiscal 2025.\footnote{The primary awardee lists this section reads are the
fiscal 2023 Live Music Fund Event Program Awardee List
(\url{https://austin.widen.net/view/pdf/jfxwxsxktx/2023LiveMusicFundEventProgramAwardeeList.pdf?t.download=true}),
the fiscal 2024 Austin Live Music Fund Awardee List
(\url{https://austin.widen.net/view/pdf/bpzuacfkk8/FY24-AwardeeList_Austin-Live-Music-Fund.pdf?t.download=true})
and the March 2026 ACME Funding Program Award List, which covers both
the fiscal 2026 Live Music Fund and the Elevate program
(\url{https://austin.widen.net/s/hhgtdnjqmn/almf_acme-2026-funding-program-award-list}).
The official program page
(\url{https://www.austintexas.gov/atxmusic/live-music-fund}) records zero
awards for fiscal 2025 and describes the cycle as deferred, with
applications folded into the fiscal 2026 round. No applicant count for
the fiscal 2026 Live Music Fund alone has been published; the city's
fiscal 2026 cultural funding news release reports 1{,}606 applicants
requesting more than \$67 million across the four current cultural
funding programs combined (Elevate, Live Music Fund, Creative Space
Assistance Program and Heritage Preservation),
\url{https://www.austintexas.gov/communications/news/city-austin-announces-2026-cultural-funding-awards},
without an individual-fund breakdown. The fiscal 2026 award list is also
inconsistent with itself in three places: it names 21 venues against its
own subtotal of 20; its musician and promoter rows fall
\NCityLmfFyTwoSixMusicianDollarGap; and its headline count of 399 awards
does not match the sum of its two section subtotals,
\NCityLmfFyTwoSixSubtotalCountGap. The share-funded column in the table
below divides press-reported award counts by press-reported applicant
counts, so the numerator differs slightly from the award count shown in
the same row (which is a tally of the city's own list PDF); the fiscal
2024 share also covers musician and promoter awards only, leaving out
that cycle's new venue awards, while the published applicant count it
divides by carries no venue-versus-nonvenue split, so the fiscal 2023 and
fiscal 2024 shares are not strictly comparable.}

\begin{center}
\footnotesize
\begin{tabular}{@{}B{2.6cm} R{2.2cm} R{3.0cm} R{3.6cm}@{}}
\toprule
\textbf{Fiscal year} & \textbf{Awards} & \textbf{Average award} &
\textbf{Share of applicants funded} \\
\midrule
2023 & \NCityLmfFyTwoThreeAwardsList{} & \NCityLmfFyTwoThreeAvg{} &
\NCityLmfFyTwoThreeFundedShare{} \\
\rowrule
2024 & \NCityLmfFyTwoFourAwards{} & \NCityLmfFyTwoFourAvg{} &
\NCityLmfFyTwoFourFundedShare{} \\
\rowrule
2025 & \multicolumn{3}{c}{No cycle ran} \\
\rowrule
2026 & \NCityLmfFyTwoSixAwards{} & \NCityLmfFyTwoSixAvg{} & Not published \\
\bottomrule
\end{tabular}
\end{center}

\vspace{3pt}
For an individual musician that variation is the operative fact: the same
application submitted in two different years faced a different chance of
funding and a different award. Predictability is part of what makes a
grant program useful for people whose income is irregular.

\subsection{Demand for venue capital ran far ahead of supply at the fund's first award}

Section~\ref{sec:venues} showed that Austin loses live-music capacity
when rooms stop reporting and no one replaces them, which points at the
cost of buildings and the security of leases. The Austin Economic
Development Corporation's Iconic Venue Fund is aimed at exactly that: it
helps venues buy or secure their buildings.\footnote{The Iconic Venue
Fund is a capital-preservation program run by the Austin Economic
Development Corporation, a public non-profit created by the city, and it
is distinct from the grantmaking Live Music Fund the rest of this
section reads. It draws on hotel-occupancy-tax revenue and a 2018 bond
package. The 2023-snapshot figures in this paragraph are from Austin
Chronicle reporting on the fund's first award: Genevieve Wood,
\textit{What Hole in the Wall's Freshly Inked 20-Year Lease Means for
Other At-Risk Venues}, 24 August 2023,
\url{https://www.austinchronicle.com/music/2023-08-25/what-hole-in-the-walls-freshly-inked-20-year-lease-means-for-other-at-risk-venues/}.
No primary Austin Economic Development Corporation document publishing
these counts was located, and this study reports the newspaper figures
as secondary.} At its first award in August 2023 the Iconic Venue Fund
had received \NCityIvfProposals\ proposals requesting
\NCityIvfRequested. Its selection committee had shortlisted
\NCityIvfShortlisted\ proposals and had made a single award of
\NCityIvfAwardTwoZeroTwoThree\ to the Hole in the Wall, so the requests
in hand ran to \NCityIvfRequestedToAwardRatio\ the size of that single
award.\footnote{The \$300 million requested figure is a stated floor
rather than an exact total, so the ratio to the single award is at least
the value shown. That figure is a 2023 snapshot rather than the fund's
current state. City Open Budget expense data record cumulative Iconic
Venue Fund grants to subrecipients of \$15.3~million through fiscal 2026
(dataset yeeq-kk6v, \texttt{fund\_nm~=~Iconic~Venue~Fund}), so
substantially more has been awarded since; against that total the \$300
million floor implies a ratio roughly an order of magnitude smaller. The
Austin Economic Development Corporation has not published an updated
proposal count, shortlist or recipient list. Cumulative Open Budget data
also put the fund's revenue minus expenditure at
\NCityIvfCumulativeNetNominalFyTwoSix\ through fiscal 2026, but that
position is produced almost entirely by fiscal 2026, an incomplete year
in which \$2.9 million of spending is recorded against revenue that has
not yet posted; summed through fiscal 2025, where the data are complete,
the fund is net positive.}

The imbalance is not confined to venue capital. In the 2020 Creative
Space Disaster Relief round, \NCityCsapTwoZeroTwoZeroApplicants\
applicants requested about \$3.7 million against the \$1 million
available, which the program page describes as nearly four times the
money on offer
(\NCityCsapTwoZeroTwoZeroOversubVsAvailable).\footnote{City of Austin,
Creative Space Assistance Program page,
\url{https://www.austintexas.gov/creative-space-assistance-program},
retrieved 1 August 2026. The 2020 round funded 32 of the 65 applicants
for \$987{,}943 total, with the awardee mix running to 12 live-music
venues, 11 arts nonprofits, 7 theaters, 1 museum or gallery and 1
individual artist. This 2020 round is the only Creative Space cycle for
which the city has published an applicant count.}

\subsection{No performance audit of the Live Music Fund exists}

The Office of the City Auditor has published no performance audit of the
Live Music Fund. Its one adjacent review, a 2024 special request on
cultural arts funding, states in its own scope that it did not look at
funding for music, and a search of the auditor's report index returns no
music-titled audit.\footnote{City of Austin Office of the City Auditor,
\textit{Special Request: EDD Culture and Arts Funding} (October 2024),
which reports on its own page 2 that ``we did not look at funding for
music or music disaster'' relief:
\url{https://austin.widen.net/view/pdf/pifyu2e2j6/Special_Request_EDD_Culture_and_Arts_Funding_October_2024.pdf?t.download=true}.
The auditor's public index of published reports at
\url{https://www.austintexas.gov/auditor/audit-reports} carries no
title matching the Live Music Fund as of 1 August 2026.} That is a
separate question from the city's annual financial audit, whose
fund-balance schedule would be needed to reconcile the Live Music Fund's
cash position against the press reporting cited in
Figure~\ref{fig:lmfflow}.

In mid-2025, 60 Live Music Fund recipients were out of compliance with
their reporting requirements, \$550,000 in awards was frozen pending
their final reports, and \$70,000 of that had been forfeited
outright.\footnote{Austin Monitor, \textit{Dozens of City Music Grants
Stalled Over Missing Final Reports}, 19 June 2025,
\url{https://austinmonitor.com/stories/2025/06/dozens-of-city-music-grants-stalled-over-missing-final-reports/},
reporting a 16 June 2025 presentation to the Arts Commission. This is
secondary reporting; no primary city document publishing these counts
was located, and no denominator for the 60 non-compliant recipients was
published alongside the count. The \$550,000 figure is the hotel-tax
share of \$847,500 frozen citywide across 70 contracts.}

The absence of a performance audit matters despite the fund's size. The
Live Music Fund is small, but its award rules have
changed in every cycle, its administration has moved to a third party,
and its cash position cannot be reconciled against public reporting
using open data alone. Those are the conditions under which a periodic
audit is normally expected.

The Live Music Fund is the city's main direct instrument for music, and it
is one of several public levers this workforce faces. Section~\ref{sec:state}
turns to the state and federal ones: the Texas Music Incubator Rebate
Program the Governor's office administers, the arts-economy aggregate the
Bureau of Economic Analysis publishes at the state level, and the federal
Shuttered Venue Operators Grant program that was the largest single
intervention of the period.
